\section{Levelwise solvability, approximation, and global obstruction}
\label{sec:E}

For a morphism $d:X\to Y$ of inverse diagrams of Hausdorff topological abelian groups, put $K_a=\ker d_a$ and $I_a=\im d_a$ as abstract groups, with their subspace topologies when topology is mentioned. Write $d_\infty:\lim X\to\lim Y$ for the induced continuous map. A \emph{levelwise solvable} element belongs to $\lim I\subseteq\lim Y$. We distinguish this subgroup from the generally larger $\lim\overline I$.

\begin{theorem}[E.1; Proved: the obstruction subgroup and the closure]\label{E:obstruction}
For every small inverse index and every such $d$, there is a natural isomorphism
\[
 \frac{\lim I}{\im d_\infty}
 \ \cong\ \ker\bigl(\lim{}^1K\longrightarrow\lim{}^1X\bigr).
\]
The class of a levelwise solution $y$ is zero if and only if $d_\infty x=y$ has a solution. In general there is no natural map from \emph{all} of $\lim^1K$ into $\coker d_\infty$; the displayed kernel is the subgroup that canonically embeds in that cokernel.

For a tower $X_0\leftarrow X_1\leftarrow\cdots$, suppose either
\begin{enumerate}[label=(\alph*)]
\item all transition maps of $X$ are surjective, with no completeness or metrizability requirement; or
\item the $X_n$ are complete metrizable groups and all transition maps have dense image.
\end{enumerate}
Then
\[
 \overline{\im d_\infty}=\lim_n\overline{I_n}
 \quad\text{inside }\lim_nY_n.
\]
In particular every levelwise solvable $y$ lies in $\overline{\im d_\infty}$. In case (a), $\lim^1X=0$, so the whole group $\lim^1K$ identifies with $\lim I/\im d_\infty$. This is the full closure quotient if every $I_n$ is closed. In general the precise additional quotient is
\[
 0\longrightarrow \frac{\lim I}{\im d_\infty}
 \longrightarrow \frac{\overline{\im d_\infty}}{\im d_\infty}
 \longrightarrow \frac{\lim\overline I}{\lim I}\longrightarrow0.
\]
The corresponding assertions hold on every directed set with a countable cofinal subset.
\end{theorem}
\begin{proof}
The objectwise exact sequence $0\to K\to X\to I\to0$ gives the initial part of the derived-limit long exact sequence
\[
 0\to\lim K\to\lim X\to\lim I
 \xrightarrow{\partial}\lim{}^1K\to\lim{}^1X.
\]
Its exactness proves the first assertion and specifies the direction of the obstruction map. Choosing local lifts $x_a$ of $y_a$ gives its usual cocycle, the differences between the transition of a lift and the chosen lift at the target.

For a tower with surjective transitions, every projection $\lim X\to X_N$ is surjective: choose a lift successively at $N+1,N+2,\ldots$ and use transitions to define the earlier coordinates. This is a set-theoretic construction and uses neither topology nor completeness.

For complete metrizable $X_n$ and dense transitions, every projection $\lim X\to X_N$ has dense image. Here is the needed approximation argument. Start with a prescribed $x_N$ and a neighborhood of it. At stage $m\geq N$ a point at level $m$ determines a finite compatible string through level $m$. Density of $X_{m+1}\to X_m$ and continuity of the finitely many earlier composites let us choose the next string so that it changes coordinate $i\leq m$ by less than $2^{-m}\varepsilon_i$, for prescribed positive error budgets. Choose the budgets at $N$ with total error inside the original neighborhood. For each fixed $i$ the resulting sequence of coordinates is Cauchy; completeness supplies a limit in $X_i$, and continuity gives compatibility. The limiting coordinate at $N$ is in the prescribed neighborhood. One may use compatible complete translation-invariant metrics, which exist for completely metrizable abelian groups; the same finite-string proof works with any specified complete compatible metrics.

Now a basic neighborhood of $y\in\lim Y$ involves finitely many coordinates, all below some $N$. If $y_N\in\overline{I_N}$, first choose $v_N\in X_N$ whose image under $d_N$ approximates $y_N$ sufficiently well in those finitely many coordinates. Surjectivity, or density, of the limit projection then supplies $v\in\lim X$ with $d_\infty v$ in the original neighborhood. This proves $\lim\overline I\subseteq\overline{\im d_\infty}$. The reverse inclusion follows because each $\overline{I_n}$ is closed and contains the coordinate image of $d_\infty$.

For a surjective tower, the endomorphism of products
\[
 \Delta_X:\prod_{n\geq0}X_n\to\prod_{n\geq0}X_n,
 \qquad (x_n)\mapsto(x_n-f_nx_{n+1})
\]
is surjective by recursion, and its cokernel computes $\lim^1X$. This proves the remaining claim in (a). The last displayed sequence is the quotient sequence for nested subgroups. If all $I_n$ are closed, its last group is zero.

For a directed set with a countable cofinal subset, recursively dominate its enumerated elements to obtain a cofinal increasing sequence. Restriction to this sequence identifies ordinary limits, with their subspace topologies, and all derived limits. For the latter, the derived cofinality comparison of Roos resolutions applies: for each object the category of sequence positions above it is a nonempty tail and hence has contractible nerve. This verifies the homological hypothesis of the comparison, rather than using ordinary cofinality alone.
\end{proof}

\begin{example}[E.2; Proved: why closure and local solvability differ]\label{E:local-closure}
Even a one-object tower with identity transitions need not identify $\lim^1K$ with the closure quotient if the images $I_n$ are not closed. Take every $X_n=\Z$ discrete, every $Y_n=\mathbb T$, and $d_n(m)=m\alpha\bmod\Z$ for an irrational $\alpha$. Then $K_n=0$, so $\lim^1K=0$, whereas
\[
 \overline{\im d_\infty}/\im d_\infty
       =\mathbb T/(\alpha\Z\bmod\Z)\ne0.
\]
All levelwise solvable compatible elements are already globally solvable. The additional closure classes fail to be levelwise solvable. This is exactly the last quotient in Theorem~\ref{E:obstruction}.
\end{example}

\begin{example}[E.3; Proved: dense transitions need completeness]\label{E:incomplete}
Let $t$ be a positive transcendental real number. For each $n\geq1$ let
\[
 X_n=\operatorname{span}_{\Q}\{(t^j,1):j\geq n\}\subset\R^2,
 \qquad Y_n=\Q\subset\R,
\]
with their subspace topologies. Use inclusions $X_{n+1}\hookrightarrow X_n$, identity transitions on $Y$, and the second-coordinate map $d_n:X_n\to\Q$. Every inclusion has dense image, because two distinct displayed vectors have a $\Q$-span dense in $\R^2$. Every $d_n$ is continuous and surjective. However $\bigcap_nX_n=0$: transcendence makes the family $(t^j,1)$ linearly independent over $\Q$, and a vector belonging to every tail span has empty finite support. Thus $\lim X=0$, while $1\in\lim Y=\Q$ is levelwise solvable and is outside the closure of the zero image. The groups $X_n$ are not complete. This refutes the version with only dense transitions and no completeness assumption.
\end{example}

\begin{theorem}[E.4; Proved: the compact regime]\label{E:compact}
For a small category $A$, the following are equivalent:
\begin{enumerate}[label=(\roman*)]
\item every levelwise solvable element for every morphism of compact Hausdorff abelian $A^{\op}$-diagrams is globally solvable;
\item $\colim_A:\Ab^A\to\Ab$ is exact.
\end{enumerate}
For every such compact morphism, $\im d_\infty$ is closed. Consequently any levelwise solvable element that is not globally solvable lies outside the closure of the image. If $A$ is directed, every positive derived limit of every compact Hausdorff abelian diagram vanishes, with no countability assumption and no assumption on the transition maps.
\end{theorem}
\begin{proof}
Images and kernels of continuous maps of compact Hausdorff groups are compact and closed. Thus $I$ is a compact group diagram and $X\to I$ is objectwise surjective. The equivalence is precisely the lifting assertion in Theorem~\ref{C:torsors}, applied to $X\to I$; conversely apply it to every objectwise surjective map. The space $\lim X$ is compact, and $\lim Y$ is Hausdorff, so its continuous image is closed. If $A$ is directed, filtered colimits of discrete abelian groups are exact. Pontryagin duality in Theorem~\ref{C:duality} gives all the asserted vanishing. This concerns the derived limits of the underlying abstract groups; it is stronger than a mere statement that the topological limit functor preserves surjections.
\end{proof}

\begin{theorem}[E.5; Refuted: surjective towers can be replaced by arbitrary directed indices]\label{E:uncountable}
There is an inverse system of complete metrizable abelian groups over the directed ordinal $\omega_1$, with surjective transitions, and an objectwise surjective morphism to the constant finite group $\F_2$, for which the source limit is zero. Thus the levelwise solvable element $1$ is outside the closure of the global image, although the index is directed and hence exact.
\end{theorem}
\begin{proof}
For each countable ordinal $\alpha$, let $S_\alpha$ be the set of injections $s:\alpha\to\N$ whose images have infinite complement. For $\alpha\leq\beta<\omega_1$, use restriction. Every restriction is surjective: partition the infinite unused set of natural numbers into two infinite sets and map the extra countably many arguments into one of them, leaving the other unused. Each $S_\alpha$ is nonempty. A compatible family would give an injection $\omega_1\to\N$, which is impossible.

Set $X_\alpha=\F_2^{(S_\alpha)}$, the free vector space with finite supports, endowed with the discrete topology. These are complete metrizable topological abelian groups, even when their underlying sets are uncountable. The induced transition maps are surjective. Let $Y_\alpha=\F_2$ be constant and let $d_\alpha$ be augmentation. All $d_\alpha$ are surjective.

We prove $\lim X=0$, including the support point that can fail for other directed sets. Suppose $(x_\alpha)$ is a coherent nonzero family. Its finite support sizes are nondecreasing with $\alpha$. They are bounded: otherwise choose countably many ordinals at which their sizes tend to infinity, and take their supremum, still below $\omega_1$. At that supremum a finite support would have to dominate every size. Thus the sizes attain a positive maximum $m$ on a tail. Every map between supports on that tail is then a bijection: all $m$ target support points must receive a nonzero coefficient, and there are only $m$ source support points, leaving no room for collisions, cancellations, or an image outside the support. Hence these finite nonempty supports form an inverse system of bijections. Choosing any point at one level and transporting it coherently gives a point of $\lim S_\alpha$, a contradiction. Therefore the source limit is zero.

The image in $\lim Y=\F_2$ is the closed zero subgroup; $1$ is levelwise solvable and has nonzero obstruction in $\lim^1\ker d_\alpha$. This is the standard support construction behind \citet[Corollary 8 and its preceding argument]{Bergman2011}, here written out for this explicit system. The countable-upper-bound property of $\omega_1$ is essential to this particular proof. We do not infer that supports automatically form a subsystem for an arbitrary directed index.
\end{proof}

\begin{proposition}[E.6; Proved: the model obstructions and their gaps]\label{E:models}
The two specified model cases have the following operators and obstruction groups.
\begin{enumerate}[label=(\alph*)]
\item On the tower indexed by positive integers, take $X_n=\Z$ with identity transitions, $Y_n=\Z/2^n\Z$ with reduction maps, and $d_n$ reduction. Then $d_\infty:\Z\to\Z_2$ has dense proper image, every finite-level equation is solvable, and
\[
 \lim{}^1(2^n\Z)\cong\Z_2/\Z.
\]
\item Over $\Cdelta^{\op}$ take $d_X=D_X^\vee:k^{P_X\,2}\to k^{E_X}$ from Proposition~\ref{C:explicit-torsor}, for $k=\F_p$, and the constant compatible right-hand side one. Every $d_X$ is surjective, but the limit map is the zero map $k^2\to k$. Its cokernel and the kernel-diagram obstruction group are both $k$. In the discrete metric on $k$, the distance of one from the limit image is one. Replacing $k$ by $\mathbb T$ and the right-hand side by $0\ne\theta$ gives the same statements with obstruction group $\mathbb T$; its distance from zero is positive for any fixed compatible metric on $\mathbb T$.
\item The system of Theorem~\ref{E:uncountable} has zero levelwise gap and limit gap one in $\F_2$, but its index is exact. Its obstruction belongs to $\ker(\lim^1K\to\lim^1X)$ and contains the injected copy of $\F_2$ supplied by the connecting map. We do not identify all of $\lim^1K$ with $\F_2$.
\end{enumerate}
\end{proposition}
\begin{proof}
For (a), the kernel tower consists of inclusions $2^{n+1}\Z\hookrightarrow2^n\Z$. The tower of source groups is surjective, so Theorem~\ref{E:obstruction} and the exact sequence of limits identify the cokernel with $\lim^1$ of those kernels. Every residue class modulo $2^n$ is represented by an integer, proving density. The image is proper since $\Z$ is countable and $\Z_2$ is uncountable, for example by its unique infinite binary expansions. With the usual $2$-adic metric the infimum distance from any $2$-adic element to the image is zero, though only integers have exact lifts.

For (b), all assertions, including nonzero cocycle representatives, were proved in Proposition~\ref{C:explicit-torsor}; representable duals have zero first derived limits, so here the full kernel obstruction group is indeed the cokernel. The assertions in (c) follow from the proved exact sequence and the explicit zero limit. This locates precisely the failure of the tower hypotheses. An obstruction outside the closure in the compact regime requires a nonexact index; outside that regime it may instead reflect failure of countable cofinality, completeness, or dense transition assumptions.
\end{proof}

\begin{theorem}[E.7; Known and independently reproved: countable derived-limit dichotomy]\label{E:gray}
Let $I$ be a directed set with a countable cofinal subset, and let $A:I^{\op}\to\Ab$ be a diagram of countable abelian groups. Then $\lim^1_I A$ is zero or uncountable. In particular a nonzero finite $\lim^1$ of a diagram of finite abelian groups cannot occur on a directed index. For arbitrary directed indices the latter assertion also follows from compact duality, since finite groups are compact.
\end{theorem}
\begin{proof}
Derived cofinality reduces to a tower, as justified in Theorem~\ref{E:obstruction}. The tower's Roos complex is the two-term product complex $\prod A_n\xrightarrow{\Delta}\prod A_n$, so $\lim^i=0$ for $i\geq2$.

Recall that the Mittag--Leffler condition says that, for every $n$, the images of $A_m\to A_n$ stabilize for large $m$. Under this condition put $B_n$ equal to that stable image. The transitions of $B$ are surjective, and the quotient tower $A/B$ is pro-zero: for each fixed $n$, every sufficiently long composite into its $n$th group is zero. Surjectivity gives $\lim^1B=0$. For a pro-zero tower the equation $\Delta x=y$ has the solution
\[
 x_n=\sum_{j\geq n}f_{n,j}y_j,
\]
where the sum is finite for each $n$. Thus $\lim^1(A/B)=0$ and the long exact sequence gives $\lim^1A=0$.

If the Mittag--Leffler condition fails, choose $m$ such that the descending groups $H_n=\im(A_{m+n}\to A_m)$ do not stabilize. Termwise surjectivity of the map from the tail tower of $A$ to $H$ and vanishing of second derived limits give a surjection $\lim^1A\twoheadrightarrow\lim^1H$. Put $G=A_m$, a countable group. The short exact sequence $0\to H\to\Delta G\to(G/H_n)\to0$ identifies
\[
 \lim{}^1H=\bigl(\lim_nG/H_n\bigr)/\im G.
\]
After passage to a subsequence there is a strict drop at each step. The surjections $G/H_{n+1}\to G/H_n$ then have at least two elements in every fiber. Recursively choosing two lifts at each level embeds the set of infinite binary sequences into $\lim G/H_n$. Thus this inverse limit is uncountable, and quotienting by the countable image of $G$ leaves an uncountable group. This proves the dichotomy.

The source was opened: \citet[p.~242, unnumbered Proposition and the parenthetical refinement following its proof]{Gray1966} proves the needed countable-system assertion and states the uncountability refinement for abelian groups. It is also explicitly stated as \citet[Proposition 3.5]{Schochet2003}. The proof above spells out the six-term-sequence refinement in the abelian case.
\end{proof}
